%

% http://aleph0.clarku.edu/~djoyce/java/elements/bookIII/bookIII.html

\subsection{Księga III} % lang-pl
\subsection{Libro III} % lang-it

\subsubsection{Definicje} % lang-pl
\subsubsection{Definizioni} % lang-it

\begin{enumerate}
    \item [3.1] Dwa okręgi są \emph{równe}, kiedy mają równe średnice (lub równoważnie, promienie). % lang-pl
    \item [3.1] Due circoli sono \emph{uguali} quando hanno diametri uguali (o equivalentemente, raggi uguali). % lang-it

    \item [3.2] Prosta jest \emph{styczna} do okręgu, jeśli ma z nim punkt wspólny, ale po przedłużeniu nie ma drugiego takiego punktu. % lang-pl
    \item [3.2] Una retta è \emph{tangente} a un circolo se ha un punto in comune con esso, ma prolungata non ha un secondo punto in comune. % lang-it

    \item [3.3] Dwa okręgi są \emph{styczne}, kiedy mają punkt wspólny, ale nie przecinają się. % lang-pl
    \item [3.3] Due circoli sono \emph{tangenti} quando hanno un punto in comune, ma non si intersecano. % lang-it

    \item [3.4] Dwie cięciwy (oryginalnie: proste) w okręgu są \emph{równoodległe} od jego środka, kiedy odcinki prostopadłe do nich opuszczone ze środka okręgu są równe. % lang-pl
    \item [3.4] Due corde (originariamente: rette) in un circolo sono \emph{equidistanti} dal suo centro quando i segmenti perpendicolari ad esse tracciati dal centro del circolo sono uguali. % lang-it

    \item [3.5] Spośród dwóch cięciw ta jest \emph{bardziej oddalona}, do której odcinek prostopadły ze środka okręgu jest dłuższy. % lang-pl
    \item [3.5] Tra due corde, quella è \emph{più lontana} dalla quale il segmento perpendicolare dal centro del circolo è più lungo. % lang-it

    \item [3.6] \emph{Odcinek kołowy} to figura ograniczona przez cięciwę okręgu (a właściwie koła) oraz jego brzeg. % lang-pl
    \item [3.6] Un \emph{segmento circolare} è una figura delimitata da una corda del circolo (o meglio, del cerchio) e dal suo bordo. % lang-it

    \item [3.7] \emph{Kątem odcinka kołowego} jest kąt między cięciwą i brzegiem koła\footnote{To sformułowanie jest dziwne, ale jest potrzebne tylko raz, w twierdzeniu 3.16}. % lang-pl
    \item [3.7] Un \emph{angolo del segmento circolare} è l'angolo tra la corda e il bordo del cerchio\footnote{Questa formulazione è strana, ma è necessaria solo una volta, nel teorema 3.16}. % lang-it

    \item [3.8] Kątem w odcinku kołowym, co nie jest (!) tym samym, co wyżej, nazywamy kąt, którego wierzchołek leży na brzegu koła, zaś ramiona przechodzą przez cięciwę, która wyznacza odcinek kołowy. % lang-pl
    \item [3.8] Un angolo nel segmento circolare, che non è (!) lo stesso di quello sopra, chiamiamo l'angolo il cui vertice si trova sul bordo del cerchio, mentre i lati passano attraverso la corda che definisce il segmento circolare. % lang-it

    \item [3.9] Kąt w odcinku kołowym wygodniej jest nazywać kątem wpisanym opartym na łuku, który łączy końce odcinka kołowego (są dwa takie łuki). % lang-pl
    \item [3.9] Un angolo nel segmento circolare è più convenientemente chiamato angolo inscritto basato sull'arco che collega le estremità del segmento circolare (ci sono due tali archi). % lang-it

    \item [3.10] Wycinek kołowy to figura ograniczona przez dwa promienie oraz łuk koła łączący te promienie. % lang-pl
    \item [3.10] Un settore circolare è una figura delimitata da due raggi e dall'arco del cerchio che collega questi raggi. % lang-it

    \item [3.11] Dwa odcinki kołowe są \emph{podobne}, kiedy mają te same kąty. % lang-pl
    \item [3.11] Due segmenti circolari sono \emph{simili} quando hanno gli stessi angoli. % lang-it
\end{enumerate}

\subsubsection{Twierdzenia} % lang-pl
\subsubsection{Teoremi} % lang-it

\begin{enumerate}
    \item [3.1] Skonstruować środek danego okręgu. % lang-pl
    \item [3.1] Costruire il centro di un dato circolo. % lang-it

    \item [3.2] Cięciwa znajduje się wewnątrz koła. % lang-pl
    \item [3.2] Una corda si trova all'interno del cerchio. % lang-it

    \item [3.3] Średnica (cięciwa przechodząca przez środek koła), która połowi inną cięciwę niebędącą średnicą, jest do niej prostopadła. Średnica, która jest prostopadła do innej cięciwy, połowi ją. % lang-pl
    \item [3.3] Il diametro (corda che passa per il centro del cerchio) che dimezza un'altra corda che non è un diametro, è perpendicolare ad essa. Il diametro che è perpendicolare a un'altra corda, la dimezza. % lang-it

    \item [3.4] Dwie cięciwy, które nie są średnicami, nie połowią się nawzajem. % lang-pl
    \item [3.4] Due corde che non sono diametri, non si dimezzano a vicenda. % lang-it

    \item [3.5] Dwa okręgi, które się przecinają, nie mogą być współśrodkowe.  % lang-pl
    \item [3.5] Due circoli che si intersecano non possono essere concentrici. % lang-it

    \item [3.6] Dwa okręgi, które są styczne, nie są koncentryczne.
    \item [3.6] Due circoli che sono tangenti non sono concentrici. % lang-it

    \item [3.7] Dane jest koło oraz punkt $F$ leżący na jego średnicy $AD$, który nie jest środkiem. Rozpatrujemy odcinki łączące punkt $F$ z innymi punktami na brzegu koła. Wtedy: % lang-pl
    \item [3.7] Dato un cerchio e un punto $F$ che si trova sul suo diametro $AD$, che non è il centro. Consideriamo i segmenti che collegano il punto $F$ con altri punti sul bordo del cerchio. Allora: % lang-it
        \begin{itemize}
        \item odcinek, który przechodzi przez środek jest najdłuższy; % lang-pl
        \item segmento che passa per il centro è il più lungo; % lang-it
        \item odcinek, który nie przechodzi przez środek, ale leży na średnicy, jest najkrótszy; % lang-pl
        \item segmento che non passa per il centro, ma si trova sul diametro, è il più corto; % lang-it
        \item odcinek, którego koniec leży dalej od wspomnianej średnicy jest dłuższy od odcinka, którego koniec leży bliżej. % lang-pl
        \item segmento il cui estremo si trova più lontano dal diametro menzionato è più lungo del segmento il cui estremo si trova più vicino. % lang-it
        \end{itemize}
    (To twierdzenie nie jest później do niczego używane.) % lang-pl
    (Questo teorema non viene più utilizzato in seguito.) % lang-it
    
    \item [3.8] Spośród wszystkich odcinków poprowadzonych z punktu leżącego na zewnątrz koła do okręgu najdłuższy jest odcinek przechodzący przez środek koła; a długość odcinków maleje wraz z oddalaniem się od tej prostej. % lang-pl
    \item [3.8] Tra tutti i segmenti tracciati da un punto situato all'esterno del cerchio al circolo, il più lungo è il segmento che passa per il centro del cerchio; e la lunghezza dei segmenti diminuisce man mano che ci si allontana da questa retta. % lang-it

    \item [3.9] Jeżeli przez punkt wewnątrz koła przechodzą trzy cięciwy równej długości, to jest środkiem koła. % lang-pl
    \item [3.9] Se tre corde di uguale lunghezza passano attraverso un punto all'interno del cerchio, allora quel punto è il centro del cerchio. % lang-it

    \item [3.10] Dwa okręgi, które się przecinają, przecinają się w dwóch punktach. % lang-pl
    \item [3.10] Due circoli che si intersecano, si intersecano in due punti. % lang-it

    \item [3.11] Prosta łącząca środki kół stycznych wewnętrznie przechodzi przez punkt styczności kół. % lang-pl
    \item [3.11] La retta che collega i centri dei cerchi tangenti internamente passa attraverso il punto di tangenza dei cerchi. % lang-it

    \item [3.12] Odcinek łączący środki kół stycznych zewnętrznie przechodzi przez punkt styczności kół. % lang-pl
    \item [3.12] Il segmento che collega i centri dei cerchi tangenti esternamente passa attraverso il punto di tangenza dei cerchi. % lang-it

    \item [3.13] Dwa różne okręgi nie mogą być styczne (wewnętrznie lub zewnętrznie) w więcej niż jednym punkcie. % lang-pl
    \item [3.13] Due circoli diversi non possono essere tangenti (internamente o esternamente) in più di un punto. % lang-it
    
    \item [3.14] Następujące warunki są równoważne: % lang-pl
    \item [3.14] Le seguenti condizioni sono equivalenti: % lang-it
        \begin{itemize}
        \item dwie cięciwy są równej długości;  % lang-pl
        \item due corde sono di uguale lunghezza;  % lang-it
        \item dwie cięciwy są równoodległe od środka koła. % lang-pl
        \item due corde sono equidistanti dal centro del cerchio. % lang-it
        \end{itemize}

    \item [3.15] Spośród dwóch cięciw ta jest dłuższa, która leży bliżej środka koła; średnica to najdłuższa cięciwa. % lang-pl
    \item [3.15] Tra due corde, quella è più lunga che si trova più vicino al centro del cerchio; il diametro è la corda più lunga. % lang-it

    \item [3.16] Prosta prostopadła do średnicy przechodząca przez jej środek nie przechodzi przez wnętrze koła. (To twierdzenie mówi też, że ,,kąt'' między prostopadłą a brzegiem koła jest mniejszy niż dowolny kąt płaski, co trudno sformalizować i nie jest później wykorzystywane). % lang-pl
    \item [3.16] Una retta perpendicolare al diametro che passa per il suo centro non passa attraverso l'interno del cerchio. (Questo teorema afferma anche che l'"angolo" tra la perpendicolare e il bordo del cerchio è minore di qualsiasi angolo piatto, il che è difficile da formalizzare e non viene più utilizzato in seguito). % lang-it

    \item [3.17] Skonstruować styczną do danego okręgu, która przechodzi przez dany punkt. % lang-pl
    \item [3.17] Costruire una tangente a un dato circolo che passa per un punto dato. % lang-it

    \item [3.18] Odcinek łączący środek koła z punktem styczności tego koła i pewnej prostej jest prostopadły do tej prostej. % lang-pl
    \item [3.18] Il segmento che collega il centro del cerchio con il punto di tangenza di quel cerchio e una certa retta è perpendicolare a quella retta. % lang-it

    \item [3.19] Prosta prostopadła do stycznej do okręgu, przechodząca przez punkt styczności, przechodzi także przez środek koła. % lang-pl
    \item [3.19] La retta perpendicolare alla tangente al cerchio, che passa per il punto di tangenza, passa anche attraverso il centro del cerchio. % lang-it

    \item [3.20] Kąt wpisany jest dwa razy mniejszy od kąta środkowego opartego na tym samym łuku. % lang-pl
    \item [3.20] L'angolo inscritto è due volte più piccolo dell'angolo al centro basato sullo stesso arco. % lang-it

    \item [3.21] Kąty wpisane oparte na tym samym łuku są równe. % lang-pl
    \item [3.21] Gli angoli inscritti basati sullo stesso arco sono uguali. % lang-it

    \item [3.22] Suma kątów przeciwległych czworokąta wpisanego w okrąg jest równa kątowi półpełnemu. % lang-pl
    \item [3.22] La somma degli angoli opposti di un quadrilatero inscritto in un cerchio è uguale a un angolo piatto. % lang-it

    \item [3.23] Na odcinku nie mona wykreślić dwóch odcinków kół podobnych, ale nie przystających do siebie. % lang-pl
    \item [3.23] Su un segmento non si possono tracciare due segmenti circolari simili, ma non congruenti tra loro. % lang-it
    % wszyscy uzywają tego samego obrazka, więc chyba nikt nit nie wie, o co tu chodzi:
    % http://aleph0.clarku.edu/~djoyce/java/elements/bookIII/propIII23.html
    % https://farside.ph.utexas.edu/Books/Euclid/Elements.pdf s. 93
    % https://www.researchgate.net/publication/403437334_Euklides_Elementy_Ksiega_III_Euclid_Elements_Book_III_Translated_Into_Polish_From_the_Text_of_Heiberg itd.
    
    \item [3.24] (wniosek z poprzedniego) Odcinki koła wykreślone na równych odcinkach są równe. % lang-pl
    \item [3.24] (conseguenza del precedente) I segmenti circolari tracciati su segmenti uguali sono uguali. % lang-it
    
    \item [3.25] Opisać koło na danym odcinku kołowym. % lang-pl
    \item [3.25] Descrivere un cerchio su un dato segmento circolare. % lang-it
    % TO TWIERDZENIE NIE JEST DALEJ UŻYWANE
    
    \item [3.26] Dwa równe kąty wpisane w dwóch równych kołach oparte są na dwóch równych łukach. % lang-pl
    \item [3.26] Gli angoli inscritti in due cerchi uguali, basati su archi uguali, sono uguali. % lang-it
    
    \item [3.27] Dwa kąty wpisane w dwa równe koła oparte na równych łukach są równe. % lang-pl
    \item [3.27] Gli angoli inscritti in due cerchi uguali, basati su archi uguali, sono uguali. % lang-it
    
    \item [3.28] Dwie cięciwy w dwóch równych kołach wyznaczają dwie pary równych łuków. % lang-pl
    \item [3.28] Le corde in due cerchi uguali raccolgono archi uguali, così che l'arco maggiore include l'arco maggiore, e l'arco minore include l'arco minore. % lang-it
    
    \item [3.29] Dwa równe łuki w dwóch równych kołach obejmują równe cięciwy. % lang-pl
    \item [3.29] Gli archi uguali in due cerchi uguali raccolgono corde uguali. % lang-it
    
    \item [3.30] Przepołowić dany łuk. % lang-pl
    \item [3.30] Dividere a metà un dato arco. % lang-it

    \item [3.31] Kąt wpisany oparty na półokręgu jest prosty, kąt wpisany oparty na krótszym odcinku jest większy od prostego, a kąt wpisany oparty na dłuższym odcinku jest mniejszy od prostego. % lang-pl
    \item [3.31] L'angolo inscritto basato su un semicerchio è retto, l'angolo inscritto basato sul segmento più corto è maggiore di un angolo retto, e l'angolo inscritto basato sul segmento più lungo è minore di un angolo retto. % lang-it
    
    \item [3.32] Twierdzenie o kącie między styczną a cięciwą. % lang-pl
    \item [3.32] Teorema sull'angolo tra la tangente e la corda. % lang-it
    
    \item [3.33] Na danym odcinku kołowym opisać odcinek koła który by zawierał kąt równy kątowi danemu. % lang-pl
    \item [3.33] Sul dato segmento circolare descrivere un arco di cerchio che contenga un angolo uguale all'angolo dato. % lang-it
    
    \item [3.34] Na danym kole wykreślić odcinek któryby zawierał kąt równy danemu kątowi. % lang-pl
    \item [3.34] Su un dato cerchio tracciare un arco che contenga un angolo uguale all'angolo dato. % lang-it
    
    \item [3.35] Twierdzenie o cięciwach (\ref{prop_intersecting_chords}). % lang-pl
    \item [3.35] Teorema sulle corde (\ref{prop_intersecting_chords}). % lang-it
    
    \item [3.36] Twierdzenie o stycznej i siecznej (\ref{twierdzenie_o_stycznej_i_siecznej}). % lang-pl
    \item [3.36] Teorema delle tangenti e delle secanti (\ref{twierdzenie_o_stycznej_i_siecznej}). % lang-it
    
    \item [3.37] Twierdzenie odwrotne do poprzedniego. % lang-pl
    \item [3.37] Teorema inverso al precedente. % lang-it
\end{enumerate}

%